%%%%%%%%%%%%%%%%%%%%%%%%%%%%%%%%%%%%%%%%%%%%%%%%%%%%%%%%%%%%
% 14.3 NUMBER SYSTEMS
% The Composition of Advanced Mathematics
%%%%%%%%%%%%%%%%%%%%%%%%%%%%%%%%%%%%%%%%%%%%%%%%%%%%%%%%%%%%

\section{Number Systems}
\label{sec:number-systems}

\prerequisite{
Basic arithmetic, algebra, set notation, and familiarity with mathematical
symbols.
}

\learningobjectives{
By the end of this reference section, the reader should be able to:
\begin{itemize}
    \item identify the major number systems used in mathematics;
    \item distinguish natural, whole, integer, rational, irrational,
          real, and complex numbers;
    \item understand the containment relationships among standard number sets;
    \item recognize field, ring, and ordered-field structure conceptually;
    \item classify numbers into the smallest standard number system containing
          them;
    \item recognize finite and infinite decimal representations;
    \item distinguish terminating and repeating decimals;
    \item understand rational representations;
    \item recognize irrational numbers;
    \item understand the real-number line;
    \item recognize algebraic and transcendental numbers;
    \item understand complex numbers and their geometric representation;
    \item compute with complex conjugates and moduli;
    \item recognize polar and exponential forms of complex numbers;
    \item understand roots of unity;
    \item recognize modular number systems;
    \item understand finite fields conceptually;
    \item recognize extended real numbers;
    \item understand cardinal-number distinctions conceptually;
    \item distinguish discrete and continuous number systems;
    \item recognize common notation conventions;
    \item avoid common classification errors.
\end{itemize}
}

%%%%%%%%%%%%%%%%%%%%%%%%%%%%%%%%%%%%%%%%%%%%%%%%%%%%%%%%%%%%
\subsection{Overview of the Standard Number Systems}
\label{subsec:number-systems-overview}
%%%%%%%%%%%%%%%%%%%%%%%%%%%%%%%%%%%%%%%%%%%%%%%%%%%%%%%%%%%%

The principal number systems used throughout elementary and advanced
mathematics are

\[
\boxed{
\mathbb{N},
\qquad
\mathbb{Z},
\qquad
\mathbb{Q},
\qquad
\mathbb{R},
\qquad
\mathbb{C}.
}
\]

They satisfy the standard containment chain

\[
\boxed{
\mathbb{N}
\subset
\mathbb{Z}
\subset
\mathbb{Q}
\subset
\mathbb{R}
\subset
\mathbb{C}.
}
\]

Each extension introduces numbers needed to solve equations that may not be
solvable in the preceding system.

%%%%%%%%%%%%%%%%%%%%%%%%%%%%%%%%%%%%%%%%%%%%%%%%%%%%%%%%%%%%
\subsection{Natural Numbers}
\label{subsec:natural-numbers-reference}
%%%%%%%%%%%%%%%%%%%%%%%%%%%%%%%%%%%%%%%%%%%%%%%%%%%%%%%%%%%%

The natural numbers are used for counting.

A common convention is

\[
\boxed{
\mathbb{N}
=
\{1,2,3,4,\ldots\}.
}
\]

Another common convention includes zero:

\[
\boxed{
\mathbb{N}
=
\{0,1,2,3,\ldots\}.
}
\]

Because both conventions occur in mathematics, a text should explicitly state
which convention it uses when zero matters.

%%%%%%%%%%%%%%%%%%%%%%%%%%%%%%%%%%%%%%%%%%%%%%%%%%%%%%%%%%%%
\subsection{Whole Numbers}
\label{subsec:whole-numbers-reference}
%%%%%%%%%%%%%%%%%%%%%%%%%%%%%%%%%%%%%%%%%%%%%%%%%%%%%%%%%%%%

The term whole numbers commonly refers to

\[
\boxed{
\{0,1,2,3,\ldots\}.
}
\]

This terminology is common in elementary mathematics but less frequently used
as a formal named number system in advanced mathematics.

%%%%%%%%%%%%%%%%%%%%%%%%%%%%%%%%%%%%%%%%%%%%%%%%%%%%%%%%%%%%
\subsection{Integers}
\label{subsec:integers-reference}
%%%%%%%%%%%%%%%%%%%%%%%%%%%%%%%%%%%%%%%%%%%%%%%%%%%%%%%%%%%%

The integers are

\[
\boxed{
\mathbb{Z}
=
\{\ldots,-3,-2,-1,0,1,2,3,\ldots\}.
}
\]

The notation \(\mathbb{Z}\) comes from the German word
\emph{Zahlen}, meaning numbers.

Integers include:

\[
\boxed{
\text{negative integers},
}
\]

\[
\boxed{
0,
}
\]

and

\[
\boxed{
\text{positive integers}.
}
\]

%%%%%%%%%%%%%%%%%%%%%%%%%%%%%%%%%%%%%%%%%%%%%%%%%%%%%%%%%%%%
\subsection{Positive and Negative Integers}
\label{subsec:positive-negative-integers-reference}
%%%%%%%%%%%%%%%%%%%%%%%%%%%%%%%%%%%%%%%%%%%%%%%%%%%%%%%%%%%%

Positive integers may be written

\[
\boxed{
\mathbb{Z}_{>0}
=
\{1,2,3,\ldots\}.
}
\]

Nonnegative integers may be written

\[
\boxed{
\mathbb{Z}_{\geq0}
=
\{0,1,2,\ldots\}.
}
\]

Negative integers are

\[
\boxed{
\mathbb{Z}_{<0}
=
\{-1,-2,-3,\ldots\}.
}
\]

%%%%%%%%%%%%%%%%%%%%%%%%%%%%%%%%%%%%%%%%%%%%%%%%%%%%%%%%%%%%
\subsection{Rational Numbers}
\label{subsec:rational-numbers-reference}
%%%%%%%%%%%%%%%%%%%%%%%%%%%%%%%%%%%%%%%%%%%%%%%%%%%%%%%%%%%%

A rational number is any number expressible as a quotient of integers:

\[
\boxed{
\mathbb{Q}
=
\left\{
\frac{a}{b}
:
a,b\in\mathbb{Z},
\ b\neq0
\right\}.
}
\]

Examples include

\[
\boxed{
\frac12,
\quad
-\frac73,
\quad
5,
\quad
0.
}
\]

Every integer is rational because

\[
\boxed{
n
=
\frac{n}{1}.
}
\]

%%%%%%%%%%%%%%%%%%%%%%%%%%%%%%%%%%%%%%%%%%%%%%%%%%%%%%%%%%%%
\subsection{Equivalent Rational Representations}
\label{subsec:equivalent-rational-representations}
%%%%%%%%%%%%%%%%%%%%%%%%%%%%%%%%%%%%%%%%%%%%%%%%%%%%%%%%%%%%

A rational number can have many fractional representations.

For example,

\[
\boxed{
\frac12
=
\frac24
=
\frac36
=
\frac{50}{100}.
}
\]

In general,

\[
\boxed{
\frac{a}{b}
=
\frac{ka}{kb}
}
\]

for every nonzero integer \(k\).

%%%%%%%%%%%%%%%%%%%%%%%%%%%%%%%%%%%%%%%%%%%%%%%%%%%%%%%%%%%%
\subsection{Terminating Decimals}
\label{subsec:terminating-decimals-reference}
%%%%%%%%%%%%%%%%%%%%%%%%%%%%%%%%%%%%%%%%%%%%%%%%%%%%%%%%%%%%

A terminating decimal is rational.

For example,

\[
\boxed{
0.125
=
\frac{125}{1000}
=
\frac18.
}
\]

Every finite decimal expansion can be expressed as an integer divided by a
power of \(10\).

%%%%%%%%%%%%%%%%%%%%%%%%%%%%%%%%%%%%%%%%%%%%%%%%%%%%%%%%%%%%
\subsection{Repeating Decimals}
\label{subsec:repeating-decimals-reference}
%%%%%%%%%%%%%%%%%%%%%%%%%%%%%%%%%%%%%%%%%%%%%%%%%%%%%%%%%%%%

Every repeating decimal is rational.

For example,

\[
x
=
0.\overline{3}.
\]

Then

\[
10x
=
3.\overline{3}.
\]

Subtracting,

\[
9x=3.
\]

Therefore,

\[
\boxed{
0.\overline{3}
=
\frac13.
}
\]

%%%%%%%%%%%%%%%%%%%%%%%%%%%%%%%%%%%%%%%%%%%%%%%%%%%%%%%%%%%%
\subsection{Worked Example: Converting a Repeating Decimal}
\label{subsec:worked-repeating-decimal}
%%%%%%%%%%%%%%%%%%%%%%%%%%%%%%%%%%%%%%%%%%%%%%%%%%%%%%%%%%%%

\begin{workedexamplebox}{Convert \(0.\overline{27}\) to a Fraction}

Let

\[
x
=
0.\overline{27}.
\]

Since two digits repeat,

\[
100x
=
27.\overline{27}.
\]

Subtract:

\[
100x-x
=
27.
\]

Thus

\[
99x=27.
\]

Therefore,

\begin{answerbox}
\[
\boxed{
x
=
\frac{27}{99}
=
\frac{3}{11}.
}
\]
\end{answerbox}

\end{workedexamplebox}

%%%%%%%%%%%%%%%%%%%%%%%%%%%%%%%%%%%%%%%%%%%%%%%%%%%%%%%%%%%%
\subsection{Decimal Characterization of Rational Numbers}
\label{subsec:rational-decimal-characterization}
%%%%%%%%%%%%%%%%%%%%%%%%%%%%%%%%%%%%%%%%%%%%%%%%%%%%%%%%%%%%

A real number is rational if and only if its decimal expansion eventually
terminates or eventually repeats.

Thus

\[
\boxed{
\text{rational}
\iff
\text{eventually periodic decimal expansion}.
}
\]

%%%%%%%%%%%%%%%%%%%%%%%%%%%%%%%%%%%%%%%%%%%%%%%%%%%%%%%%%%%%
\subsection{Irrational Numbers}
\label{subsec:irrational-numbers-reference}
%%%%%%%%%%%%%%%%%%%%%%%%%%%%%%%%%%%%%%%%%%%%%%%%%%%%%%%%%%%%

An irrational number is a real number that is not rational.

Thus

\[
\boxed{
\mathbb{R}\setminus\mathbb{Q}
}
\]

is the set of irrational numbers.

Examples include

\[
\boxed{
\sqrt2,
\quad
\sqrt3,
\quad
\pi,
\quad
e.
}
\]

Their decimal expansions are nonterminating and nonrepeating.

%%%%%%%%%%%%%%%%%%%%%%%%%%%%%%%%%%%%%%%%%%%%%%%%%%%%%%%%%%%%
\subsection{Worked Example: Proving \(\sqrt2\) Is Irrational}
\label{subsec:worked-sqrt2-irrational}
%%%%%%%%%%%%%%%%%%%%%%%%%%%%%%%%%%%%%%%%%%%%%%%%%%%%%%%%%%%%

\begin{workedexamplebox}{A Classical Contradiction Proof}

Assume

\[
\sqrt2
=
\frac{a}{b}
\]

with integers \(a,b\) having no common factor and \(b\neq0\).

Then

\[
2
=
\frac{a^2}{b^2},
\]

so

\[
a^2
=
2b^2.
\]

Therefore \(a^2\) is even, so \(a\) is even.

Write

\[
a=2k.
\]

Then

\[
4k^2
=
2b^2,
\]

so

\[
b^2
=
2k^2.
\]

Thus \(b\) is also even.

This contradicts the assumption that \(a/b\) was in lowest terms.

\begin{answerbox}
\[
\boxed{
\sqrt2\notin\mathbb{Q}.
}
\]
\end{answerbox}

\end{workedexamplebox}

%%%%%%%%%%%%%%%%%%%%%%%%%%%%%%%%%%%%%%%%%%%%%%%%%%%%%%%%%%%%
\subsection{Real Numbers}
\label{subsec:real-numbers-reference}
%%%%%%%%%%%%%%%%%%%%%%%%%%%%%%%%%%%%%%%%%%%%%%%%%%%%%%%%%%%%

The real numbers are denoted

\[
\boxed{
\mathbb{R}.
}
\]

They consist of rational and irrational numbers:

\[
\boxed{
\mathbb{R}
=
\mathbb{Q}
\cup
(\mathbb{R}\setminus\mathbb{Q}).
}
\]

Every point on the ordinary number line corresponds to a real number.

%%%%%%%%%%%%%%%%%%%%%%%%%%%%%%%%%%%%%%%%%%%%%%%%%%%%%%%%%%%%
\subsection{The Real Number Line}
\label{subsec:real-number-line-reference}
%%%%%%%%%%%%%%%%%%%%%%%%%%%%%%%%%%%%%%%%%%%%%%%%%%%%%%%%%%%%

The real numbers possess a natural order.

If

\[
a<b,
\]

then \(a\) lies to the left of \(b\) on the real line.

The number line extends without bound:

\[
\boxed{
-\infty
\longleftarrow
\mathbb{R}
\longrightarrow
+\infty.
}
\]

The symbols \(\pm\infty\) are not ordinary real numbers.

%%%%%%%%%%%%%%%%%%%%%%%%%%%%%%%%%%%%%%%%%%%%%%%%%%%%%%%%%%%%
\subsection{Ordered Field Structure of the Reals}
\label{subsec:real-ordered-field}
%%%%%%%%%%%%%%%%%%%%%%%%%%%%%%%%%%%%%%%%%%%%%%%%%%%%%%%%%%%%

The real numbers form an ordered field.

This means one can perform

\[
\boxed{
+,\ -,\ \times,\ \div
}
\]

with division by nonzero numbers, while maintaining a compatible ordering.

For example,

\[
a<b
\]

implies

\[
\boxed{
a+c<b+c.
}
\]

If \(c>0\), then

\[
\boxed{
ac<bc.
}
\]

%%%%%%%%%%%%%%%%%%%%%%%%%%%%%%%%%%%%%%%%%%%%%%%%%%%%%%%%%%%%
\subsection{Completeness of the Real Numbers}
\label{subsec:real-completeness-reference}
%%%%%%%%%%%%%%%%%%%%%%%%%%%%%%%%%%%%%%%%%%%%%%%%%%%%%%%%%%%%

One of the most important properties distinguishing \(\mathbb{R}\) from
\(\mathbb{Q}\) is completeness.

A standard form of the completeness axiom states:

\begin{importantbox}
Every nonempty subset of \(\mathbb{R}\) that is bounded above has a least
upper bound in \(\mathbb{R}\).
\end{importantbox}

This least upper bound is the supremum.

Completeness underlies much of real analysis.

%%%%%%%%%%%%%%%%%%%%%%%%%%%%%%%%%%%%%%%%%%%%%%%%%%%%%%%%%%%%
\subsection{Density of the Rational Numbers}
\label{subsec:rational-density-reference}
%%%%%%%%%%%%%%%%%%%%%%%%%%%%%%%%%%%%%%%%%%%%%%%%%%%%%%%%%%%%

Between any two distinct real numbers there exists a rational number.

If

\[
a<b,
\]

then there exists

\[
q\in\mathbb{Q}
\]

such that

\[
\boxed{
a<q<b.
}
\]

Thus \(\mathbb{Q}\) is dense in \(\mathbb{R}\).

%%%%%%%%%%%%%%%%%%%%%%%%%%%%%%%%%%%%%%%%%%%%%%%%%%%%%%%%%%%%
\subsection{Density of the Irrational Numbers}
\label{subsec:irrational-density-reference}
%%%%%%%%%%%%%%%%%%%%%%%%%%%%%%%%%%%%%%%%%%%%%%%%%%%%%%%%%%%%

The irrational numbers are also dense in the real line.

Between any two distinct real numbers there exists an irrational number.

Therefore every real interval contains both rational and irrational numbers.

%%%%%%%%%%%%%%%%%%%%%%%%%%%%%%%%%%%%%%%%%%%%%%%%%%%%%%%%%%%%
\subsection{Algebraic Numbers}
\label{subsec:algebraic-numbers-reference}
%%%%%%%%%%%%%%%%%%%%%%%%%%%%%%%%%%%%%%%%%%%%%%%%%%%%%%%%%%%%

A complex number \(\alpha\) is algebraic if it is a root of a nonzero
polynomial with integer coefficients.

Thus

\[
\boxed{
a_n\alpha^n
+
a_{n-1}\alpha^{n-1}
+
\cdots
+
a_0
=
0
}
\]

for some integers \(a_i\), not all zero.

Examples include

\[
\boxed{
\sqrt2,
\quad
\sqrt[3]{5},
\quad
i.
}
\]

%%%%%%%%%%%%%%%%%%%%%%%%%%%%%%%%%%%%%%%%%%%%%%%%%%%%%%%%%%%%
\subsection{Transcendental Numbers}
\label{subsec:transcendental-numbers-reference}
%%%%%%%%%%%%%%%%%%%%%%%%%%%%%%%%%%%%%%%%%%%%%%%%%%%%%%%%%%%%

A number that is not algebraic is called transcendental.

Important examples include

\[
\boxed{
\pi
}
\]

and

\[
\boxed{
e.
}
\]

Thus transcendental numbers form a proper subset of the irrational numbers.

%%%%%%%%%%%%%%%%%%%%%%%%%%%%%%%%%%%%%%%%%%%%%%%%%%%%%%%%%%%%
\subsection{Algebraic Versus Irrational}
\label{subsec:algebraic-versus-irrational}
%%%%%%%%%%%%%%%%%%%%%%%%%%%%%%%%%%%%%%%%%%%%%%%%%%%%%%%%%%%%

Do not confuse the classifications.

For example,

\[
\boxed{
\sqrt2
}
\]

is irrational but algebraic because it satisfies

\[
x^2-2=0.
\]

By contrast,

\[
\boxed{
\pi
}
\]

is irrational and transcendental.

%%%%%%%%%%%%%%%%%%%%%%%%%%%%%%%%%%%%%%%%%%%%%%%%%%%%%%%%%%%%
\subsection{Complex Numbers}
\label{subsec:complex-numbers-reference}
%%%%%%%%%%%%%%%%%%%%%%%%%%%%%%%%%%%%%%%%%%%%%%%%%%%%%%%%%%%%

The complex numbers are

\[
\boxed{
\mathbb{C}
=
\{
a+bi
:
a,b\in\mathbb{R}
\},
}
\]

where

\[
\boxed{
i^2=-1.
}
\]

Every real number is complex because

\[
\boxed{
a
=
a+0i.
}
\]

Therefore,

\[
\boxed{
\mathbb{R}
\subset
\mathbb{C}.
}
\]

%%%%%%%%%%%%%%%%%%%%%%%%%%%%%%%%%%%%%%%%%%%%%%%%%%%%%%%%%%%%
\subsection{Real and Imaginary Parts}
\label{subsec:real-imaginary-parts-reference}
%%%%%%%%%%%%%%%%%%%%%%%%%%%%%%%%%%%%%%%%%%%%%%%%%%%%%%%%%%%%

For

\[
z=a+bi,
\]

the real part is

\[
\boxed{
\operatorname{Re}(z)=a,
}
\]

and the imaginary part is

\[
\boxed{
\operatorname{Im}(z)=b.
}
\]

A purely imaginary number has real part zero.

%%%%%%%%%%%%%%%%%%%%%%%%%%%%%%%%%%%%%%%%%%%%%%%%%%%%%%%%%%%%
\subsection{Complex Conjugate}
\label{subsec:complex-conjugate-number-systems}
%%%%%%%%%%%%%%%%%%%%%%%%%%%%%%%%%%%%%%%%%%%%%%%%%%%%%%%%%%%%

For

\[
z=a+bi,
\]

the complex conjugate is

\[
\boxed{
\overline{z}
=
a-bi.
}
\]

Then

\[
\boxed{
z\overline{z}
=
a^2+b^2.
}
\]

%%%%%%%%%%%%%%%%%%%%%%%%%%%%%%%%%%%%%%%%%%%%%%%%%%%%%%%%%%%%
\subsection{Complex Modulus}
\label{subsec:complex-modulus-reference}
%%%%%%%%%%%%%%%%%%%%%%%%%%%%%%%%%%%%%%%%%%%%%%%%%%%%%%%%%%%%

The modulus of

\[
z=a+bi
\]

is

\[
\boxed{
|z|
=
\sqrt{a^2+b^2}.
}
\]

Thus

\[
\boxed{
|z|^2
=
z\overline{z}.
}
\]

Geometrically, \(|z|\) is the distance from \(z\) to the origin in the
complex plane.

%%%%%%%%%%%%%%%%%%%%%%%%%%%%%%%%%%%%%%%%%%%%%%%%%%%%%%%%%%%%
\subsection{Worked Example: Complex Modulus}
\label{subsec:worked-complex-modulus}
%%%%%%%%%%%%%%%%%%%%%%%%%%%%%%%%%%%%%%%%%%%%%%%%%%%%%%%%%%%%

\begin{workedexamplebox}{Find the Modulus}

Let

\[
z=3-4i.
\]

Then

\[
|z|
=
\sqrt{
3^2+(-4)^2
}.
\]

Therefore,

\begin{answerbox}
\[
\boxed{
|z|=5.
}
\]
\end{answerbox}

\end{workedexamplebox}

%%%%%%%%%%%%%%%%%%%%%%%%%%%%%%%%%%%%%%%%%%%%%%%%%%%%%%%%%%%%
\subsection{Complex Division}
\label{subsec:complex-division-reference}
%%%%%%%%%%%%%%%%%%%%%%%%%%%%%%%%%%%%%%%%%%%%%%%%%%%%%%%%%%%%

To divide complex numbers, multiply numerator and denominator by the conjugate
of the denominator.

For

\[
\frac{a+bi}{c+di},
\]

multiply by

\[
\frac{c-di}{c-di}.
\]

Then

\[
\boxed{
\frac{a+bi}{c+di}
=
\frac{
(a+bi)(c-di)
}{
c^2+d^2
}.
}
\]

%%%%%%%%%%%%%%%%%%%%%%%%%%%%%%%%%%%%%%%%%%%%%%%%%%%%%%%%%%%%
\subsection{Complex Plane}
\label{subsec:complex-plane-reference}
%%%%%%%%%%%%%%%%%%%%%%%%%%%%%%%%%%%%%%%%%%%%%%%%%%%%%%%%%%%%

A complex number

\[
z=a+bi
\]

corresponds to the point

\[
\boxed{
(a,b)
}
\]

in the complex plane.

The horizontal axis is the real axis.

The vertical axis is the imaginary axis.

%%%%%%%%%%%%%%%%%%%%%%%%%%%%%%%%%%%%%%%%%%%%%%%%%%%%%%%%%%%%
\subsection{Polar Form}
\label{subsec:complex-polar-form-reference}
%%%%%%%%%%%%%%%%%%%%%%%%%%%%%%%%%%%%%%%%%%%%%%%%%%%%%%%%%%%%

A nonzero complex number can be written

\[
\boxed{
z
=
r(\cos\theta+i\sin\theta),
}
\]

where

\[
r=|z|
\]

and

\[
\theta
\]

is an argument of \(z\).

%%%%%%%%%%%%%%%%%%%%%%%%%%%%%%%%%%%%%%%%%%%%%%%%%%%%%%%%%%%%
\subsection{Euler Form}
\label{subsec:complex-euler-form-reference}
%%%%%%%%%%%%%%%%%%%%%%%%%%%%%%%%%%%%%%%%%%%%%%%%%%%%%%%%%%%%

Euler's identity

\[
\boxed{
e^{i\theta}
=
\cos\theta+i\sin\theta
}
\]

allows complex numbers to be written

\[
\boxed{
z
=
re^{i\theta}.
}
\]

This form simplifies multiplication, division, powers, and roots.

%%%%%%%%%%%%%%%%%%%%%%%%%%%%%%%%%%%%%%%%%%%%%%%%%%%%%%%%%%%%
\subsection{De Moivre's Formula}
\label{subsec:de-moivre-reference}
%%%%%%%%%%%%%%%%%%%%%%%%%%%%%%%%%%%%%%%%%%%%%%%%%%%%%%%%%%%%

For integer \(n\),

\[
\boxed{
(\cos\theta+i\sin\theta)^n
=
\cos(n\theta)
+
i\sin(n\theta).
}
\]

Equivalently,

\[
\boxed{
(re^{i\theta})^n
=
r^ne^{in\theta}.
}
\]

%%%%%%%%%%%%%%%%%%%%%%%%%%%%%%%%%%%%%%%%%%%%%%%%%%%%%%%%%%%%
\subsection{Roots of Complex Numbers}
\label{subsec:complex-roots-reference}
%%%%%%%%%%%%%%%%%%%%%%%%%%%%%%%%%%%%%%%%%%%%%%%%%%%%%%%%%%%%

To solve

\[
z^n
=
re^{i\theta},
\]

the \(n\) roots are

\[
\boxed{
z_k
=
r^{1/n}
e^{
i(\theta+2\pi k)/n
},
\qquad
k=0,1,\ldots,n-1.
}
\]

Thus a nonzero complex number has exactly \(n\) distinct \(n\)-th roots.

%%%%%%%%%%%%%%%%%%%%%%%%%%%%%%%%%%%%%%%%%%%%%%%%%%%%%%%%%%%%
\subsection{Roots of Unity}
\label{subsec:roots-of-unity-reference}
%%%%%%%%%%%%%%%%%%%%%%%%%%%%%%%%%%%%%%%%%%%%%%%%%%%%%%%%%%%%

The \(n\)-th roots of unity solve

\[
\boxed{
z^n=1.
}
\]

They are

\[
\boxed{
z_k
=
e^{2\pi i k/n},
\qquad
k=0,\ldots,n-1.
}
\]

These points are equally spaced around the unit circle.

%%%%%%%%%%%%%%%%%%%%%%%%%%%%%%%%%%%%%%%%%%%%%%%%%%%%%%%%%%%%
\subsection{Worked Example: Fourth Roots of Unity}
\label{subsec:worked-fourth-roots-unity}
%%%%%%%%%%%%%%%%%%%%%%%%%%%%%%%%%%%%%%%%%%%%%%%%%%%%%%%%%%%%

\begin{workedexamplebox}{Solve \(z^4=1\)}

The roots are

\[
z_k
=
e^{2\pi i k/4}.
\]

Thus

\[
z_0=1,
\]

\[
z_1=i,
\]

\[
z_2=-1,
\]

\[
z_3=-i.
\]

Therefore,

\begin{answerbox}
\[
\boxed{
\{1,i,-1,-i\}.
}
\]
\end{answerbox}

\end{workedexamplebox}

%%%%%%%%%%%%%%%%%%%%%%%%%%%%%%%%%%%%%%%%%%%%%%%%%%%%%%%%%%%%
\subsection{Fundamental Theorem of Algebra}
\label{subsec:fundamental-theorem-algebra-reference}
%%%%%%%%%%%%%%%%%%%%%%%%%%%%%%%%%%%%%%%%%%%%%%%%%%%%%%%%%%%%

A central reason for working over \(\mathbb{C}\) is the Fundamental Theorem
of Algebra.

\begin{importantbox}
Every nonconstant polynomial with complex coefficients has at least one
complex root.
\end{importantbox}

Equivalently, a polynomial of degree \(n\geq1\) factors into \(n\) linear
factors over \(\mathbb{C}\), counting multiplicity.

%%%%%%%%%%%%%%%%%%%%%%%%%%%%%%%%%%%%%%%%%%%%%%%%%%%%%%%%%%%%
\subsection{Why Number Systems Expand}
\label{subsec:why-number-systems-expand}
%%%%%%%%%%%%%%%%%%%%%%%%%%%%%%%%%%%%%%%%%%%%%%%%%%%%%%%%%%%%

Different equations motivate different number systems.

For example,

\[
\boxed{
x+3=1
}
\]

requires negative integers.

The equation

\[
\boxed{
2x=1
}
\]

requires rational numbers.

The equation

\[
\boxed{
x^2=2
}
\]

requires irrational real numbers.

The equation

\[
\boxed{
x^2=-1
}
\]

requires complex numbers.

Thus

\[
\boxed{
\mathbb{N}
\rightarrow
\mathbb{Z}
\rightarrow
\mathbb{Q}
\rightarrow
\mathbb{R}
\rightarrow
\mathbb{C}
}
\]

can be viewed as a sequence of algebraic extensions motivated by solvability.

%%%%%%%%%%%%%%%%%%%%%%%%%%%%%%%%%%%%%%%%%%%%%%%%%%%%%%%%%%%%
\subsection{Closure Properties}
\label{subsec:number-system-closure}
%%%%%%%%%%%%%%%%%%%%%%%%%%%%%%%%%%%%%%%%%%%%%%%%%%%%%%%%%%%%

A set is closed under an operation if applying that operation to members of
the set always produces another member.

For example,

\[
\boxed{
\mathbb{Z}
}
\]

is closed under addition, subtraction, and multiplication.

However, it is not closed under division because

\[
\boxed{
1\div2
=
\frac12
\notin
\mathbb{Z}.
}
\]

%%%%%%%%%%%%%%%%%%%%%%%%%%%%%%%%%%%%%%%%%%%%%%%%%%%%%%%%%%%%
\subsection{Closure Table}
\label{subsec:number-system-closure-table}
%%%%%%%%%%%%%%%%%%%%%%%%%%%%%%%%%%%%%%%%%%%%%%%%%%%%%%%%%%%%

\begin{center}
\begin{tabular}{c|c|c|c|c}
\textbf{Set}
&
\textbf{Addition}
&
\textbf{Subtraction}
&
\textbf{Multiplication}
&
\textbf{Division by Nonzero}
\\ \hline
\(\mathbb{N}\) & Yes & No & Yes & No \\
\(\mathbb{Z}\) & Yes & Yes & Yes & No \\
\(\mathbb{Q}\) & Yes & Yes & Yes & Yes \\
\(\mathbb{R}\) & Yes & Yes & Yes & Yes \\
\(\mathbb{C}\) & Yes & Yes & Yes & Yes
\end{tabular}
\end{center}

Here the natural-number convention is assumed not to alter the relevant
closure conclusions.

%%%%%%%%%%%%%%%%%%%%%%%%%%%%%%%%%%%%%%%%%%%%%%%%%%%%%%%%%%%%
\subsection{Fields}
\label{subsec:number-fields-reference}
%%%%%%%%%%%%%%%%%%%%%%%%%%%%%%%%%%%%%%%%%%%%%%%%%%%%%%%%%%%%

The systems

\[
\boxed{
\mathbb{Q},
\qquad
\mathbb{R},
\qquad
\mathbb{C}
}
\]

are fields.

Informally, a field permits:

\[
\boxed{
\text{addition},
}
\]

\[
\boxed{
\text{subtraction},
}
\]

\[
\boxed{
\text{multiplication},
}
\]

and

\[
\boxed{
\text{division by nonzero elements}.
}
\]

The integers do not form a field.

%%%%%%%%%%%%%%%%%%%%%%%%%%%%%%%%%%%%%%%%%%%%%%%%%%%%%%%%%%%%
\subsection{Ordered Fields}
\label{subsec:ordered-fields-reference}
%%%%%%%%%%%%%%%%%%%%%%%%%%%%%%%%%%%%%%%%%%%%%%%%%%%%%%%%%%%%

The rational and real numbers are ordered fields.

Thus

\[
\boxed{
\mathbb{Q}
}
\]

and

\[
\boxed{
\mathbb{R}
}
\]

have compatible arithmetic and order structures.

The complex numbers cannot be ordered in a way that makes them an ordered
field extending the usual real-number ordering.

%%%%%%%%%%%%%%%%%%%%%%%%%%%%%%%%%%%%%%%%%%%%%%%%%%%%%%%%%%%%
\subsection{Finite Number Systems Modulo \(n\)}
\label{subsec:modular-number-systems}
%%%%%%%%%%%%%%%%%%%%%%%%%%%%%%%%%%%%%%%%%%%%%%%%%%%%%%%%%%%%

For positive integer \(n\), the residue classes modulo \(n\) form

\[
\boxed{
\mathbb{Z}/n\mathbb{Z}.
}
\]

This set may be represented by

\[
\boxed{
\{0,1,\ldots,n-1\}
}
\]

with arithmetic performed modulo \(n\).

For example,

\[
\boxed{
7+5
\equiv
2
\pmod{10}.
}
\]

%%%%%%%%%%%%%%%%%%%%%%%%%%%%%%%%%%%%%%%%%%%%%%%%%%%%%%%%%%%%
\subsection{Arithmetic Modulo \(n\)}
\label{subsec:modular-arithmetic-reference}
%%%%%%%%%%%%%%%%%%%%%%%%%%%%%%%%%%%%%%%%%%%%%%%%%%%%%%%%%%%%

Two integers are congruent modulo \(n\) when

\[
\boxed{
a\equiv b\pmod n
}
\]

if and only if

\[
\boxed{
n\mid(a-b).
}
\]

For example,

\[
\boxed{
17\equiv5\pmod{12}.
}
\]

%%%%%%%%%%%%%%%%%%%%%%%%%%%%%%%%%%%%%%%%%%%%%%%%%%%%%%%%%%%%
\subsection{Finite Fields}
\label{subsec:finite-fields-reference}
%%%%%%%%%%%%%%%%%%%%%%%%%%%%%%%%%%%%%%%%%%%%%%%%%%%%%%%%%%%%

When \(p\) is prime,

\[
\boxed{
\mathbb{Z}/p\mathbb{Z}
}
\]

is a field.

It is often denoted

\[
\boxed{
\mathbb{F}_p.
}
\]

For example,

\[
\boxed{
\mathbb{F}_5
=
\{0,1,2,3,4\}
}
\]

with arithmetic modulo \(5\).

%%%%%%%%%%%%%%%%%%%%%%%%%%%%%%%%%%%%%%%%%%%%%%%%%%%%%%%%%%%%
\subsection{Worked Example: Inverses Modulo a Prime}
\label{subsec:worked-modular-inverse-prime}
%%%%%%%%%%%%%%%%%%%%%%%%%%%%%%%%%%%%%%%%%%%%%%%%%%%%%%%%%%%%

\begin{workedexamplebox}{Find \(3^{-1}\) in \(\mathbb{F}_7\)}

We seek \(x\) such that

\[
3x
\equiv
1
\pmod7.
\]

Testing residues,

\[
3(5)
=
15
\equiv
1
\pmod7.
\]

Therefore,

\begin{answerbox}
\[
\boxed{
3^{-1}
\equiv
5
\pmod7.
}
\]
\end{answerbox}

\end{workedexamplebox}

%%%%%%%%%%%%%%%%%%%%%%%%%%%%%%%%%%%%%%%%%%%%%%%%%%%%%%%%%%%%
\subsection{Prime Fields}
\label{subsec:prime-fields-reference}
%%%%%%%%%%%%%%%%%%%%%%%%%%%%%%%%%%%%%%%%%%%%%%%%%%%%%%%%%%%%

The simplest finite fields have prime order:

\[
\boxed{
\mathbb{F}_p.
}
\]

More generally, finite fields exist with

\[
\boxed{
p^n
}
\]

elements where \(p\) is prime and \(n\geq1\).

Up to isomorphism, there is exactly one finite field of each prime-power
order.

%%%%%%%%%%%%%%%%%%%%%%%%%%%%%%%%%%%%%%%%%%%%%%%%%%%%%%%%%%%%
\subsection{Extended Real Numbers}
\label{subsec:extended-real-numbers-reference}
%%%%%%%%%%%%%%%%%%%%%%%%%%%%%%%%%%%%%%%%%%%%%%%%%%%%%%%%%%%%

In analysis, one sometimes adjoins positive and negative infinity to the real
numbers:

\[
\boxed{
\overline{\mathbb{R}}
=
\mathbb{R}
\cup
\{-\infty,+\infty\}.
}
\]

This is called the extended real number system.

The symbols \(\pm\infty\) are not ordinary real numbers.

%%%%%%%%%%%%%%%%%%%%%%%%%%%%%%%%%%%%%%%%%%%%%%%%%%%%%%%%%%%%
\subsection{Use of the Extended Reals}
\label{subsec:extended-real-use}
%%%%%%%%%%%%%%%%%%%%%%%%%%%%%%%%%%%%%%%%%%%%%%%%%%%%%%%%%%%%

Extended reals are useful for expressions such as

\[
\boxed{
\sup A
=
+\infty
}
\]

or extended-valued functions

\[
\boxed{
f:X
\to
(-\infty,+\infty].
}
\]

Such notation appears frequently in optimization and measure theory.

%%%%%%%%%%%%%%%%%%%%%%%%%%%%%%%%%%%%%%%%%%%%%%%%%%%%%%%%%%%%
\subsection{Projective Extension of the Complex Numbers}
\label{subsec:riemann-sphere-reference}
%%%%%%%%%%%%%%%%%%%%%%%%%%%%%%%%%%%%%%%%%%%%%%%%%%%%%%%%%%%%

Complex analysis sometimes adds one point at infinity:

\[
\boxed{
\widehat{\mathbb{C}}
=
\mathbb{C}
\cup
\{\infty\}.
}
\]

This is the extended complex plane or Riemann sphere.

It differs from the extended real line, which distinguishes \(+\infty\) and
\(-\infty\).

%%%%%%%%%%%%%%%%%%%%%%%%%%%%%%%%%%%%%%%%%%%%%%%%%%%%%%%%%%%%
\subsection{Gaussian Integers}
\label{subsec:gaussian-integers-reference}
%%%%%%%%%%%%%%%%%%%%%%%%%%%%%%%%%%%%%%%%%%%%%%%%%%%%%%%%%%%%

The Gaussian integers are

\[
\boxed{
\mathbb{Z}[i]
=
\{
a+bi
:
a,b\in\mathbb{Z}
\}.
}
\]

They form a subring of \(\mathbb{C}\).

Examples include

\[
\boxed{
3+2i,
\quad
-1+5i,
\quad
7.
}
\]

%%%%%%%%%%%%%%%%%%%%%%%%%%%%%%%%%%%%%%%%%%%%%%%%%%%%%%%%%%%%
\subsection{Quadratic Number Extensions}
\label{subsec:quadratic-number-extensions-reference}
%%%%%%%%%%%%%%%%%%%%%%%%%%%%%%%%%%%%%%%%%%%%%%%%%%%%%%%%%%%%

A quadratic extension of \(\mathbb{Q}\) may have form

\[
\boxed{
\mathbb{Q}(\sqrt d)
=
\{
a+b\sqrt d
:
a,b\in\mathbb{Q}
\},
}
\]

for suitable squarefree integer \(d\).

For example,

\[
\boxed{
\mathbb{Q}(\sqrt2).
}
\]

These structures appear in algebraic number theory.

%%%%%%%%%%%%%%%%%%%%%%%%%%%%%%%%%%%%%%%%%%%%%%%%%%%%%%%%%%%%
\subsection{Real Algebraic Numbers}
\label{subsec:real-algebraic-numbers}
%%%%%%%%%%%%%%%%%%%%%%%%%%%%%%%%%%%%%%%%%%%%%%%%%%%%%%%%%%%%

The set of real algebraic numbers contains:

\[
\boxed{
\mathbb{Q}
}
\]

and many irrational numbers such as

\[
\boxed{
\sqrt2,
\quad
\sqrt[3]{7}.
}
\]

It forms a field contained in \(\mathbb{R}\).

%%%%%%%%%%%%%%%%%%%%%%%%%%%%%%%%%%%%%%%%%%%%%%%%%%%%%%%%%%%%
\subsection{Cardinality of the Natural Numbers}
\label{subsec:natural-number-cardinality}
%%%%%%%%%%%%%%%%%%%%%%%%%%%%%%%%%%%%%%%%%%%%%%%%%%%%%%%%%%%%

The natural numbers are countably infinite.

Their cardinality is denoted

\[
\boxed{
\aleph_0.
}
\]

A set is countable if it is finite or can be placed in one-to-one
correspondence with \(\mathbb{N}\).

%%%%%%%%%%%%%%%%%%%%%%%%%%%%%%%%%%%%%%%%%%%%%%%%%%%%%%%%%%%%
\subsection{Countability of the Integers}
\label{subsec:integer-countability-reference}
%%%%%%%%%%%%%%%%%%%%%%%%%%%%%%%%%%%%%%%%%%%%%%%%%%%%%%%%%%%%

Although the integers extend infinitely in both directions,

\[
\boxed{
|\mathbb{Z}|
=
|\mathbb{N}|
=
\aleph_0.
}
\]

Thus \(\mathbb{Z}\) is countably infinite.

%%%%%%%%%%%%%%%%%%%%%%%%%%%%%%%%%%%%%%%%%%%%%%%%%%%%%%%%%%%%
\subsection{Countability of the Rational Numbers}
\label{subsec:rational-countability-reference}
%%%%%%%%%%%%%%%%%%%%%%%%%%%%%%%%%%%%%%%%%%%%%%%%%%%%%%%%%%%%

The rational numbers are also countable:

\[
\boxed{
|\mathbb{Q}|
=
\aleph_0.
}
\]

This may seem surprising because rationals are dense in \(\mathbb{R}\).

Density and cardinality are different concepts.

%%%%%%%%%%%%%%%%%%%%%%%%%%%%%%%%%%%%%%%%%%%%%%%%%%%%%%%%%%%%
\subsection{Uncountability of the Real Numbers}
\label{subsec:real-uncountability-reference}
%%%%%%%%%%%%%%%%%%%%%%%%%%%%%%%%%%%%%%%%%%%%%%%%%%%%%%%%%%%%

The real numbers are uncountable:

\[
\boxed{
|\mathbb{R}|
>
|\mathbb{N}|.
}
\]

Cantor's diagonal argument proves that no enumeration of all real numbers is
possible.

The cardinality of \(\mathbb{R}\) is often called the continuum.

%%%%%%%%%%%%%%%%%%%%%%%%%%%%%%%%%%%%%%%%%%%%%%%%%%%%%%%%%%%%
\subsection{Cardinality of the Complex Numbers}
\label{subsec:complex-cardinality-reference}
%%%%%%%%%%%%%%%%%%%%%%%%%%%%%%%%%%%%%%%%%%%%%%%%%%%%%%%%%%%%

Because

\[
\mathbb{C}
\cong
\mathbb{R}^2
\]

as sets,

\[
\boxed{
|\mathbb{C}|
=
|\mathbb{R}|.
}
\]

Thus the complex plane has the same cardinality as the real line.

%%%%%%%%%%%%%%%%%%%%%%%%%%%%%%%%%%%%%%%%%%%%%%%%%%%%%%%%%%%%
\subsection{Countable Versus Dense}
\label{subsec:countable-versus-dense-reference}
%%%%%%%%%%%%%%%%%%%%%%%%%%%%%%%%%%%%%%%%%%%%%%%%%%%%%%%%%%%%

A set can be countable and dense.

For example,

\[
\boxed{
\mathbb{Q}
}
\]

is countable but dense in \(\mathbb{R}\).

Therefore,

\[
\boxed{
\text{dense}
\not\Rightarrow
\text{uncountable}.
}
\]

%%%%%%%%%%%%%%%%%%%%%%%%%%%%%%%%%%%%%%%%%%%%%%%%%%%%%%%%%%%%
\subsection{Discrete and Continuous Number Sets}
\label{subsec:discrete-continuous-number-sets}
%%%%%%%%%%%%%%%%%%%%%%%%%%%%%%%%%%%%%%%%%%%%%%%%%%%%%%%%%%%%

The integers are discrete because each integer is isolated from the others.

The real numbers form a continuum.

Informally,

\[
\boxed{
\mathbb{Z}
=
\text{discrete},
}
\]

while

\[
\boxed{
\mathbb{R}
=
\text{continuous}.
}
\]

This distinction is important in optimization, probability, and dynamical
systems.

%%%%%%%%%%%%%%%%%%%%%%%%%%%%%%%%%%%%%%%%%%%%%%%%%%%%%%%%%%%%
\subsection{Intervals of Real Numbers}
\label{subsec:real-intervals-number-systems}
%%%%%%%%%%%%%%%%%%%%%%%%%%%%%%%%%%%%%%%%%%%%%%%%%%%%%%%%%%%%

Common real intervals include

\[
\boxed{
(a,b),
\quad
[a,b],
\quad
(a,b],
\quad
[a,b).
}
\]

Unbounded intervals include

\[
\boxed{
(-\infty,a),
}
\]

\[
\boxed{
[a,\infty).
}
\]

Infinity is never included as a real endpoint.

%%%%%%%%%%%%%%%%%%%%%%%%%%%%%%%%%%%%%%%%%%%%%%%%%%%%%%%%%%%%
\subsection{Positive Real Numbers}
\label{subsec:positive-real-numbers-reference}
%%%%%%%%%%%%%%%%%%%%%%%%%%%%%%%%%%%%%%%%%%%%%%%%%%%%%%%%%%%%

The positive real numbers may be written

\[
\boxed{
\mathbb{R}_{>0}
}
\]

or sometimes

\[
\boxed{
\mathbb{R}_+.
}
\]

Because \(\mathbb{R}_+\) is used inconsistently across texts, the notation

\[
\boxed{
\mathbb{R}_{>0}
}
\]

is less ambiguous when zero must be excluded.

%%%%%%%%%%%%%%%%%%%%%%%%%%%%%%%%%%%%%%%%%%%%%%%%%%%%%%%%%%%%
\subsection{Nonnegative Real Numbers}
\label{subsec:nonnegative-real-reference}
%%%%%%%%%%%%%%%%%%%%%%%%%%%%%%%%%%%%%%%%%%%%%%%%%%%%%%%%%%%%

The nonnegative real numbers are

\[
\boxed{
\mathbb{R}_{\geq0}
=
[0,\infty).
}
\]

Similarly,

\[
\boxed{
\mathbb{R}_{\leq0}
=
(-\infty,0].
}
\]

%%%%%%%%%%%%%%%%%%%%%%%%%%%%%%%%%%%%%%%%%%%%%%%%%%%%%%%%%%%%
\subsection{Real Vector Spaces}
\label{subsec:real-vector-space-reference}
%%%%%%%%%%%%%%%%%%%%%%%%%%%%%%%%%%%%%%%%%%%%%%%%%%%%%%%%%%%%

The notation

\[
\boxed{
\mathbb{R}^n
}
\]

means the set of ordered \(n\)-tuples of real numbers:

\[
\boxed{
\mathbb{R}^n
=
\{
(x_1,\ldots,x_n)
:
x_i\in\mathbb{R}
\}.
}
\]

For example,

\[
\mathbb{R}^2
\]

is the real coordinate plane.

%%%%%%%%%%%%%%%%%%%%%%%%%%%%%%%%%%%%%%%%%%%%%%%%%%%%%%%%%%%%
\subsection{Complex Vector Spaces}
\label{subsec:complex-vector-space-reference}
%%%%%%%%%%%%%%%%%%%%%%%%%%%%%%%%%%%%%%%%%%%%%%%%%%%%%%%%%%%%

Similarly,

\[
\boxed{
\mathbb{C}^n
}
\]

is the set of ordered \(n\)-tuples of complex numbers.

It is fundamental in complex linear algebra, Fourier analysis, quantum
mechanics, and spectral theory.

%%%%%%%%%%%%%%%%%%%%%%%%%%%%%%%%%%%%%%%%%%%%%%%%%%%%%%%%%%%%
\subsection{Matrices Over Number Systems}
\label{subsec:matrices-over-number-systems}
%%%%%%%%%%%%%%%%%%%%%%%%%%%%%%%%%%%%%%%%%%%%%%%%%%%%%%%%%%%%

The notation

\[
\boxed{
\mathbb{R}^{m\times n}
}
\]

means the set of \(m\times n\) real matrices.

Likewise,

\[
\boxed{
\mathbb{C}^{m\times n}
}
\]

denotes complex matrices.

More generally,

\[
\boxed{
\mathbb{F}^{m\times n}
}
\]

may denote matrices over a field \(\mathbb{F}\).

%%%%%%%%%%%%%%%%%%%%%%%%%%%%%%%%%%%%%%%%%%%%%%%%%%%%%%%%%%%%
\subsection{Classification Hierarchy}
\label{subsec:number-classification-hierarchy}
%%%%%%%%%%%%%%%%%%%%%%%%%%%%%%%%%%%%%%%%%%%%%%%%%%%%%%%%%%%%

For standard elementary classification, use the hierarchy

\[
\boxed{
\mathbb{N}
\subset
\mathbb{Z}
\subset
\mathbb{Q}
\subset
\mathbb{R}
\subset
\mathbb{C}.
}
\]

For example,

\[
5
\]

belongs to every set in the chain.

However, its smallest standard classification is

\[
\boxed{
5\in\mathbb{N}.
}
\]

%%%%%%%%%%%%%%%%%%%%%%%%%%%%%%%%%%%%%%%%%%%%%%%%%%%%%%%%%%%%
\subsection{Worked Example: Classifying Numbers}
\label{subsec:worked-classifying-numbers}
%%%%%%%%%%%%%%%%%%%%%%%%%%%%%%%%%%%%%%%%%%%%%%%%%%%%%%%%%%%%

\begin{workedexamplebox}{Smallest Standard Number Set}

Classify each number using the smallest standard set.

\[
-7,
\qquad
\frac{3}{5},
\qquad
\sqrt2,
\qquad
4i.
\]

For \(-7\),

\[
-7\in\mathbb{Z}.
\]

For \(3/5\),

\[
\frac35\in\mathbb{Q}.
\]

For \(\sqrt2\),

\[
\sqrt2\in\mathbb{R}\setminus\mathbb{Q}.
\]

For \(4i\),

\[
4i\in\mathbb{C}\setminus\mathbb{R}.
\]

\begin{answerbox}
\[
\boxed{
-7:\mathbb{Z},
\qquad
\frac35:\mathbb{Q},
\qquad
\sqrt2:\text{irrational real},
\qquad
4i:\mathbb{C}.
}
\]
\end{answerbox}

\end{workedexamplebox}

%%%%%%%%%%%%%%%%%%%%%%%%%%%%%%%%%%%%%%%%%%%%%%%%%%%%%%%%%%%%
\subsection{Radicals and Number Classification}
\label{subsec:radicals-number-classification}
%%%%%%%%%%%%%%%%%%%%%%%%%%%%%%%%%%%%%%%%%%%%%%%%%%%%%%%%%%%%

A radical is not automatically irrational.

For example,

\[
\boxed{
\sqrt9=3
}
\]

is natural.

Similarly,

\[
\boxed{
\sqrt{\frac14}
=
\frac12
}
\]

is rational.

By contrast,

\[
\boxed{
\sqrt2
}
\]

is irrational.

The value must be simplified before classification.

%%%%%%%%%%%%%%%%%%%%%%%%%%%%%%%%%%%%%%%%%%%%%%%%%%%%%%%%%%%%
\subsection{Negative Square Roots}
\label{subsec:negative-square-root-classification}
%%%%%%%%%%%%%%%%%%%%%%%%%%%%%%%%%%%%%%%%%%%%%%%%%%%%%%%%%%%%

Over the real numbers,

\[
\sqrt{-1}
\]

is not defined as a real number.

In the complex system,

\[
\boxed{
\sqrt{-1}=i.
}
\]

Thus

\[
\boxed{
\sqrt{-9}=3i
}
\]

using the principal square root convention.

%%%%%%%%%%%%%%%%%%%%%%%%%%%%%%%%%%%%%%%%%%%%%%%%%%%%%%%%%%%%
\subsection{Fractions and Rationality}
\label{subsec:fractions-rationality-reference}
%%%%%%%%%%%%%%%%%%%%%%%%%%%%%%%%%%%%%%%%%%%%%%%%%%%%%%%%%%%%

A fraction is rational only when numerator and denominator are integers with
nonzero denominator.

For example,

\[
\boxed{
\frac{\sqrt2}{2}
}
\]

is not rational simply because it is written as a quotient.

Since \(\sqrt2\) is irrational,

\[
\boxed{
\frac{\sqrt2}{2}
}
\]

is irrational.

%%%%%%%%%%%%%%%%%%%%%%%%%%%%%%%%%%%%%%%%%%%%%%%%%%%%%%%%%%%%
\subsection{Decimal Approximation Versus Exact Value}
\label{subsec:decimal-approximation-exact-value}
%%%%%%%%%%%%%%%%%%%%%%%%%%%%%%%%%%%%%%%%%%%%%%%%%%%%%%%%%%%%

An irrational number can be approximated by rational decimals.

For example,

\[
\boxed{
\pi
\approx
3.14159.
}
\]

But

\[
\boxed{
\pi
\neq
3.14159.
}
\]

The decimal \(3.14159\) is rational because it terminates.

%%%%%%%%%%%%%%%%%%%%%%%%%%%%%%%%%%%%%%%%%%%%%%%%%%%%%%%%%%%%
\subsection{Scientific Notation}
\label{subsec:scientific-notation-reference}
%%%%%%%%%%%%%%%%%%%%%%%%%%%%%%%%%%%%%%%%%%%%%%%%%%%%%%%%%%%%

A nonzero real number may be written in scientific notation as

\[
\boxed{
a\times10^n,
}
\]

where typically

\[
\boxed{
1\leq|a|<10.
}
\]

For example,

\[
\boxed{
420000
=
4.2\times10^5.
}
\]

%%%%%%%%%%%%%%%%%%%%%%%%%%%%%%%%%%%%%%%%%%%%%%%%%%%%%%%%%%%%
\subsection{Binary Numbers}
\label{subsec:binary-numbers-reference}
%%%%%%%%%%%%%%%%%%%%%%%%%%%%%%%%%%%%%%%%%%%%%%%%%%%%%%%%%%%%

Base \(2\) uses digits

\[
\boxed{
0,\ 1.
}
\]

For example,

\[
\boxed{
1011_2
}
\]

means

\[
1(2^3)
+
0(2^2)
+
1(2^1)
+
1(2^0).
\]

Thus

\[
\boxed{
1011_2
=
11_{10}.
}
\]

%%%%%%%%%%%%%%%%%%%%%%%%%%%%%%%%%%%%%%%%%%%%%%%%%%%%%%%%%%%%
\subsection{General Positional Number Systems}
\label{subsec:positional-number-systems}
%%%%%%%%%%%%%%%%%%%%%%%%%%%%%%%%%%%%%%%%%%%%%%%%%%%%%%%%%%%%

In base \(b\),

\[
\boxed{
a_na_{n-1}\cdots a_1a_0
}
\]

represents

\[
\boxed{
\sum_{k=0}^{n}
a_kb^k,
}
\]

where

\[
0\leq a_k<b.
\]

Common bases include:

\[
\boxed{
2,\ 8,\ 10,\ 16.
}
\]

%%%%%%%%%%%%%%%%%%%%%%%%%%%%%%%%%%%%%%%%%%%%%%%%%%%%%%%%%%%%
\subsection{Hexadecimal Numbers}
\label{subsec:hexadecimal-reference}
%%%%%%%%%%%%%%%%%%%%%%%%%%%%%%%%%%%%%%%%%%%%%%%%%%%%%%%%%%%%

Base \(16\) uses symbols

\[
\boxed{
0,1,\ldots,9,A,B,C,D,E,F.
}
\]

The letters represent

\[
\boxed{
A=10,
\quad
B=11,
\quad
\ldots,
\quad
F=15.
}
\]

For example,

\[
\boxed{
1F_{16}
=
1(16)+15
=
31_{10}.
}
\]

%%%%%%%%%%%%%%%%%%%%%%%%%%%%%%%%%%%%%%%%%%%%%%%%%%%%%%%%%%%%
\subsection{Number Systems and Computation}
\label{subsec:number-systems-computation}
%%%%%%%%%%%%%%%%%%%%%%%%%%%%%%%%%%%%%%%%%%%%%%%%%%%%%%%%%%%%

Computers cannot generally store arbitrary real numbers exactly.

Floating-point arithmetic represents only a finite subset of rational numbers.

Thus a computed value such as

\[
\boxed{
0.1
}
\]

may not have an exact finite binary representation.

This creates rounding error.

%%%%%%%%%%%%%%%%%%%%%%%%%%%%%%%%%%%%%%%%%%%%%%%%%%%%%%%%%%%%
\subsection{Floating-Point Numbers}
\label{subsec:floating-point-reference}
%%%%%%%%%%%%%%%%%%%%%%%%%%%%%%%%%%%%%%%%%%%%%%%%%%%%%%%%%%%%

A simplified floating-point representation has form

\[
\boxed{
\pm m\times b^e,
}
\]

where:

\[
m
=
\text{significand},
\]

\[
b
=
\text{base},
\]

and

\[
e
=
\text{exponent}.
\]

Floating-point numbers approximate real arithmetic but do not reproduce it
exactly.

%%%%%%%%%%%%%%%%%%%%%%%%%%%%%%%%%%%%%%%%%%%%%%%%%%%%%%%%%%%%
\subsection{Machine Precision}
\label{subsec:machine-precision-reference}
%%%%%%%%%%%%%%%%%%%%%%%%%%%%%%%%%%%%%%%%%%%%%%%%%%%%%%%%%%%%

Machine precision describes the scale at which floating-point numbers can
distinguish nearby values around a given magnitude.

Consequently, in numerical computation,

\[
\boxed{
x+y=x
}
\]

can occur numerically when \(y\) is extremely small relative to \(x\).

This is a computational phenomenon, not a failure of real arithmetic.

%%%%%%%%%%%%%%%%%%%%%%%%%%%%%%%%%%%%%%%%%%%%%%%%%%%%%%%%%%%%
\subsection{Special Number Subsets}
\label{subsec:special-number-subsets}
%%%%%%%%%%%%%%%%%%%%%%%%%%%%%%%%%%%%%%%%%%%%%%%%%%%%%%%%%%%%

Several important subsets of the integers occur frequently.

These include:

\[
\boxed{
\text{even integers},
}
\]

\[
\boxed{
\text{odd integers},
}
\]

\[
\boxed{
\text{prime numbers},
}
\]

and

\[
\boxed{
\text{composite numbers}.
}
\]

These are classifications inside \(\mathbb{Z}\) or \(\mathbb{N}\), not new
number systems comparable with \(\mathbb{R}\) or \(\mathbb{C}\).

%%%%%%%%%%%%%%%%%%%%%%%%%%%%%%%%%%%%%%%%%%%%%%%%%%%%%%%%%%%%
\subsection{Even Integers}
\label{subsec:even-integers-reference}
%%%%%%%%%%%%%%%%%%%%%%%%%%%%%%%%%%%%%%%%%%%%%%%%%%%%%%%%%%%%

An integer \(n\) is even if

\[
\boxed{
n=2k
}
\]

for some

\[
k\in\mathbb{Z}.
\]

Examples include

\[
\boxed{
-4,
\quad
0,
\quad
12.
}
\]

%%%%%%%%%%%%%%%%%%%%%%%%%%%%%%%%%%%%%%%%%%%%%%%%%%%%%%%%%%%%
\subsection{Odd Integers}
\label{subsec:odd-integers-reference}
%%%%%%%%%%%%%%%%%%%%%%%%%%%%%%%%%%%%%%%%%%%%%%%%%%%%%%%%%%%%

An integer \(n\) is odd if

\[
\boxed{
n=2k+1
}
\]

for some

\[
k\in\mathbb{Z}.
\]

Every integer is exactly one of even or odd.

%%%%%%%%%%%%%%%%%%%%%%%%%%%%%%%%%%%%%%%%%%%%%%%%%%%%%%%%%%%%
\subsection{Prime Numbers}
\label{subsec:prime-numbers-reference}
%%%%%%%%%%%%%%%%%%%%%%%%%%%%%%%%%%%%%%%%%%%%%%%%%%%%%%%%%%%%

A prime number is an integer

\[
p>1
\]

whose only positive divisors are

\[
\boxed{
1
\quad\text{and}\quad
p.
}
\]

Examples include

\[
\boxed{
2,3,5,7,11,13,\ldots
}
\]

The number \(1\) is not prime.

%%%%%%%%%%%%%%%%%%%%%%%%%%%%%%%%%%%%%%%%%%%%%%%%%%%%%%%%%%%%
\subsection{Composite Numbers}
\label{subsec:composite-numbers-reference}
%%%%%%%%%%%%%%%%%%%%%%%%%%%%%%%%%%%%%%%%%%%%%%%%%%%%%%%%%%%%

A positive integer greater than \(1\) that is not prime is composite.

Examples include

\[
\boxed{
4,
\quad
6,
\quad
8,
\quad
9,
\quad
10.
}
\]

%%%%%%%%%%%%%%%%%%%%%%%%%%%%%%%%%%%%%%%%%%%%%%%%%%%%%%%%%%%%
\subsection{Fundamental Theorem of Arithmetic}
\label{subsec:fundamental-theorem-arithmetic-reference}
%%%%%%%%%%%%%%%%%%%%%%%%%%%%%%%%%%%%%%%%%%%%%%%%%%%%%%%%%%%%

Every integer \(n>1\) can be expressed as a product of primes, uniquely up to
the order of the factors.

For example,

\[
\boxed{
60
=
2^2\cdot3\cdot5.
}
\]

This prime factorization structure is fundamental to number theory.

%%%%%%%%%%%%%%%%%%%%%%%%%%%%%%%%%%%%%%%%%%%%%%%%%%%%%%%%%%%%
\subsection{Quick Classification Table}
\label{subsec:number-systems-classification-table}
%%%%%%%%%%%%%%%%%%%%%%%%%%%%%%%%%%%%%%%%%%%%%%%%%%%%%%%%%%%%

\begin{center}
\begin{tabular}{c|l|l}
\textbf{System}
&
\textbf{Notation}
&
\textbf{Description}
\\ \hline
Natural numbers
&
\(\mathbb{N}\)
&
counting numbers, convention may include \(0\)
\\

Integers
&
\(\mathbb{Z}\)
&
positive and negative whole numbers and \(0\)
\\

Rational numbers
&
\(\mathbb{Q}\)
&
ratios \(a/b\) of integers
\\

Irrational numbers
&
\(\mathbb{R}\setminus\mathbb{Q}\)
&
real but not rational
\\

Real numbers
&
\(\mathbb{R}\)
&
all points on the real line
\\

Complex numbers
&
\(\mathbb{C}\)
&
numbers \(a+bi\)
\\

Algebraic numbers
&
context dependent
&
roots of nonzero integer polynomials
\\

Transcendental numbers
&
context dependent
&
numbers that are not algebraic
\end{tabular}
\end{center}

%%%%%%%%%%%%%%%%%%%%%%%%%%%%%%%%%%%%%%%%%%%%%%%%%%%%%%%%%%%%
\subsection{Quick Examples}
\label{subsec:number-systems-quick-examples}
%%%%%%%%%%%%%%%%%%%%%%%%%%%%%%%%%%%%%%%%%%%%%%%%%%%%%%%%%%%%

\begin{center}
\begin{tabular}{c|l}
\textbf{Number} & \textbf{Smallest Common Classification} \\ \hline
\(7\) & natural \\
\(0\) & whole/nonnegative integer \\
\(-8\) & integer \\
\(\frac25\) & rational \\
\(0.125\) & rational \\
\(0.\overline{12}\) & rational \\
\(\sqrt2\) & irrational real \\
\(\pi\) & transcendental real \\
\(e\) & transcendental real \\
\(3+2i\) & complex \\
\(5i\) & purely imaginary complex
\end{tabular}
\end{center}

%%%%%%%%%%%%%%%%%%%%%%%%%%%%%%%%%%%%%%%%%%%%%%%%%%%%%%%%%%%%
\subsection{Common Errors}
\label{subsec:number-systems-common-errors}
%%%%%%%%%%%%%%%%%%%%%%%%%%%%%%%%%%%%%%%%%%%%%%%%%%%%%%%%%%%%

\subsubsection{Assuming Natural Numbers Always Include Zero}

Both conventions occur:

\[
\mathbb{N}
=
\{1,2,3,\ldots\}
\]

and

\[
\mathbb{N}
=
\{0,1,2,3,\ldots\}.
\]

The convention should be stated.

%%%%%%%%%%%%%%%%%%%%%%%%%%%%%%%%%%%%%%%%%%%%%%%%%%%%%%%%%%%%
\subsubsection{Calling Every Decimal Irrational}

A decimal is rational if it terminates or eventually repeats.

For example,

\[
\boxed{
0.25
=
\frac14.
}
\]

%%%%%%%%%%%%%%%%%%%%%%%%%%%%%%%%%%%%%%%%%%%%%%%%%%%%%%%%%%%%
\subsubsection{Calling Every Nonterminating Decimal Irrational}

Repeating decimals are nonterminating but rational.

For example,

\[
\boxed{
0.\overline{6}
=
\frac23.
}
\]

%%%%%%%%%%%%%%%%%%%%%%%%%%%%%%%%%%%%%%%%%%%%%%%%%%%%%%%%%%%%
\subsubsection{Assuming Every Radical Is Irrational}

For example,

\[
\boxed{
\sqrt{16}
=
4.
}
\]

Always simplify first.

%%%%%%%%%%%%%%%%%%%%%%%%%%%%%%%%%%%%%%%%%%%%%%%%%%%%%%%%%%%%
\subsubsection{Assuming Every Irrational Number Is Transcendental}

The number

\[
\boxed{
\sqrt2
}
\]

is irrational but algebraic.

Transcendental is a stronger classification.

%%%%%%%%%%%%%%%%%%%%%%%%%%%%%%%%%%%%%%%%%%%%%%%%%%%%%%%%%%%%
\subsubsection{Assuming Fractions Are Automatically Rational}

A rational number must be a ratio of integers.

For example,

\[
\boxed{
\frac{\pi}{2}
}
\]

is not rational.

%%%%%%%%%%%%%%%%%%%%%%%%%%%%%%%%%%%%%%%%%%%%%%%%%%%%%%%%%%%%
\subsubsection{Assuming \(i\) Is a Real Number}

The number

\[
\boxed{
i
}
\]

is complex and satisfies

\[
i^2=-1.
\]

It is not real.

%%%%%%%%%%%%%%%%%%%%%%%%%%%%%%%%%%%%%%%%%%%%%%%%%%%%%%%%%%%%
\subsubsection{Treating Infinity as a Real Number}

The symbol

\[
\boxed{
\infty
}
\]

does not represent an ordinary real number.

For example,

\[
\boxed{
\infty+1=\infty
}
\]

is not ordinary real-number arithmetic.

Its interpretation depends on an extended mathematical context.

%%%%%%%%%%%%%%%%%%%%%%%%%%%%%%%%%%%%%%%%%%%%%%%%%%%%%%%%%%%%
\subsubsection{Assuming \(\mathbb{Q}\) Is Sparse Because It Is Countable}

The rational numbers are countable but dense in \(\mathbb{R}\).

Every real interval contains infinitely many rational numbers.

%%%%%%%%%%%%%%%%%%%%%%%%%%%%%%%%%%%%%%%%%%%%%%%%%%%%%%%%%%%%
\subsubsection{Assuming \(\mathbb{R}\) and \(\mathbb{C}\) Have Different Cardinalities}

Although \(\mathbb{C}\) is geometrically two-dimensional,

\[
\boxed{
|\mathbb{C}|
=
|\mathbb{R}|.
}
\]

Dimension and cardinality are different concepts.

%%%%%%%%%%%%%%%%%%%%%%%%%%%%%%%%%%%%%%%%%%%%%%%%%%%%%%%%%%%%
\subsubsection{Confusing Modular Arithmetic With Ordinary Equality}

The statement

\[
\boxed{
7\equiv2\pmod5
}
\]

does not mean

\[
7=2.
\]

It means their difference is divisible by \(5\).

%%%%%%%%%%%%%%%%%%%%%%%%%%%%%%%%%%%%%%%%%%%%%%%%%%%%%%%%%%%%
\subsubsection{Assuming \(\mathbb{Z}/n\mathbb{Z}\) Is Always a Field}

It is a field exactly when \(n\) is prime.

For example,

\[
\boxed{
\mathbb{Z}/5\mathbb{Z}
}
\]

is a field, but

\[
\boxed{
\mathbb{Z}/6\mathbb{Z}
}
\]

is not.

%%%%%%%%%%%%%%%%%%%%%%%%%%%%%%%%%%%%%%%%%%%%%%%%%%%%%%%%%%%%
\subsubsection{Confusing Exact Numbers With Floating-Point Approximations}

A computer representation of

\[
\boxed{
\sqrt2
}
\]

is normally only an approximation.

The exact mathematical number and its stored decimal approximation are not
the same object.

%%%%%%%%%%%%%%%%%%%%%%%%%%%%%%%%%%%%%%%%%%%%%%%%%%%%%%%%%%%%
\subsection{Exercises}
\label{subsec:number-systems-reference-exercises}
%%%%%%%%%%%%%%%%%%%%%%%%%%%%%%%%%%%%%%%%%%%%%%%%%%%%%%%%%%%%

\begin{exercise}[1]
List the standard containment chain from \(\mathbb{N}\) through
\(\mathbb{C}\).
\end{exercise}

\begin{exercise}[1]
Define a rational number.
\end{exercise}

\begin{exercise}[1]
Define an irrational number.
\end{exercise}

\begin{exercise}[1]
Define a complex number.
\end{exercise}

\begin{exercise}[1]
State one convention for the natural numbers.
\end{exercise}

\begin{exercise}[2]
Classify

\[
-12
\]

using the smallest standard number system.
\end{exercise}

\begin{exercise}[2]
Classify

\[
\frac{17}{4}.
\]
\end{exercise}

\begin{exercise}[2]
Classify

\[
\sqrt{49}.
\]
\end{exercise}

\begin{exercise}[2]
Classify

\[
\sqrt7.
\]
\end{exercise}

\begin{exercise}[2]
Classify

\[
2-5i.
\]
\end{exercise}

\begin{exercise}[2]
Convert

\[
0.\overline{6}
\]

to a fraction.
\end{exercise}

\begin{exercise}[2]
Convert

\[
0.\overline{18}
\]

to a fraction.
\end{exercise}

\begin{exercise}[2]
Explain why every terminating decimal is rational.
\end{exercise}

\begin{exercise}[2]
Explain why every repeating decimal is rational.
\end{exercise}

\begin{exercise}[3]
Prove that

\[
\sqrt3
\]

is irrational.
\end{exercise}

\begin{exercise}[2]
Determine whether

\[
\sqrt{25}
\]

is rational.
\end{exercise}

\begin{exercise}[2]
Determine whether

\[
\frac{\sqrt2}{4}
\]

is rational.
\end{exercise}

\begin{exercise}[2]
Find the real and imaginary parts of

\[
z=-3+7i.
\]
\end{exercise}

\begin{exercise}[2]
Find the conjugate of

\[
z=4-9i.
\]
\end{exercise}

\begin{exercise}[2]
Compute

\[
|5+12i|.
\]
\end{exercise}

\begin{exercise}[2]
Divide

\[
\frac{3+2i}{1-i}.
\]
\end{exercise}

\begin{exercise}[3]
Convert a complex number from rectangular to polar form.
\end{exercise}

\begin{exercise}[3]
Use De Moivre's formula to compute

\[
(1+i)^8.
\]
\end{exercise}

\begin{exercise}[3]
Find all cube roots of unity.
\end{exercise}

\begin{exercise}[3]
Find all fourth roots of

\[
16.
\]

Interpret the roots in \(\mathbb{C}\).
\end{exercise}

\begin{exercise}[2]
Explain why \(\mathbb{Z}\) is not closed under division.
\end{exercise}

\begin{exercise}[2]
Explain why \(\mathbb{Q}\) is a field.
\end{exercise}

\begin{exercise}[3]
Explain why \(\mathbb{C}\) cannot be made into an ordered field extending the
usual order on \(\mathbb{R}\).
\end{exercise}

\begin{exercise}[2]
Compute

\[
13+11
\pmod 7.
\]
\end{exercise}

\begin{exercise}[2]
Compute

\[
6\cdot5
\pmod 7.
\]
\end{exercise}

\begin{exercise}[3]
Find the multiplicative inverse of \(4\) modulo \(7\).
\end{exercise}

\begin{exercise}[3]
Explain why every nonzero element of \(\mathbb{F}_p\) has an inverse.
\end{exercise}

\begin{exercise}[2]
Determine whether

\[
\mathbb{Z}/8\mathbb{Z}
\]

is a field.
\end{exercise}

\begin{exercise}[2]
Define an algebraic number.
\end{exercise}

\begin{exercise}[2]
Define a transcendental number.
\end{exercise}

\begin{exercise}[2]
Classify \(\sqrt5\) as algebraic or transcendental.
\end{exercise}

\begin{exercise}[2]
State whether \(\pi\) is algebraic or transcendental.
\end{exercise}

\begin{exercise}[3]
Explain why every rational number is algebraic.
\end{exercise}

\begin{exercise}[2]
Define the extended real numbers.
\end{exercise}

\begin{exercise}[2]
Explain why \(+\infty\) is not an ordinary real number.
\end{exercise}

\begin{exercise}[2]
Define the Gaussian integers.
\end{exercise}

\begin{exercise}[3]
Determine whether

\[
3+4i
\]

is a Gaussian integer.
\end{exercise}

\begin{exercise}[2]
Explain what it means for a set to be countable.
\end{exercise}

\begin{exercise}[2]
State whether \(\mathbb{Z}\) is countable.
\end{exercise}

\begin{exercise}[2]
State whether \(\mathbb{Q}\) is countable.
\end{exercise}

\begin{exercise}[2]
State whether \(\mathbb{R}\) is countable.
\end{exercise}

\begin{exercise}[3]
Explain how \(\mathbb{Q}\) can be both countable and dense.
\end{exercise}

\begin{exercise}[3]
Explain why

\[
|\mathbb{C}|=|\mathbb{R}|.
\]
\end{exercise}

\begin{exercise}[2]
Convert

\[
101101_2
\]

to base \(10\).
\end{exercise}

\begin{exercise}[2]
Convert

\[
45_{10}
\]

to binary.
\end{exercise}

\begin{exercise}[2]
Convert

\[
2A_{16}
\]

to base \(10\).
\end{exercise}

\begin{exercise}[3]
Explain why \(0.1\) cannot generally be represented exactly by a finite
binary floating-point expansion.
\end{exercise}

\begin{exercise}[3]
Compare exact real arithmetic with floating-point arithmetic.
\end{exercise}

\begin{exercise}[2]
Define an even integer.
\end{exercise}

\begin{exercise}[2]
Define an odd integer.
\end{exercise}

\begin{exercise}[2]
Define a prime number.
\end{exercise}

\begin{exercise}[2]
Explain why \(1\) is not prime.
\end{exercise}

\begin{exercise}[2]
Factor

\[
840
\]

into primes.
\end{exercise}

\begin{exercise}[3]
State the Fundamental Theorem of Arithmetic.
\end{exercise}

\begin{exercise}[4]
Construct an explicit enumeration showing that \(\mathbb{Z}\) is countable.
\end{exercise}

\begin{exercise}[4]
Outline an argument showing that \(\mathbb{Q}\) is countable.
\end{exercise}

\begin{exercise}[4]
Study Cantor's diagonal argument proving that \(\mathbb{R}\) is uncountable.
\end{exercise}

\begin{exercise}[4]
Prove that the algebraic numbers are countable.
\end{exercise}

\begin{exercise}[4]
Use the previous result and the uncountability of \(\mathbb{R}\) to show that
transcendental real numbers exist.
\end{exercise}

\begin{exercise}[4]
Study the field structure of

\[
\mathbb{Q}(\sqrt2).
\]
\end{exercise}

\begin{exercise}[4]
Construct the finite field \(\mathbb{F}_4\) conceptually.
\end{exercise}

\begin{exercise}[4]
Study the relationship between roots of unity and cyclic groups.
\end{exercise}

\begin{exercise}[4]
Study the Riemann sphere as the extended complex plane.
\end{exercise}

%%%%%%%%%%%%%%%%%%%%%%%%%%%%%%%%%%%%%%%%%%%%%%%%%%%%%%%%%%%%
\subsection{Conceptual Summary}
\label{subsec:number-systems-summary}
%%%%%%%%%%%%%%%%%%%%%%%%%%%%%%%%%%%%%%%%%%%%%%%%%%%%%%%%%%%%

The principal number systems satisfy

\[
\boxed{
\mathbb{N}
\subset
\mathbb{Z}
\subset
\mathbb{Q}
\subset
\mathbb{R}
\subset
\mathbb{C}.
}
\]

Natural numbers support counting.

Integers add negative quantities.

Rational numbers allow quotients.

Real numbers include irrational quantities and form a complete ordered field.

Complex numbers introduce

\[
\boxed{
i^2=-1
}
\]

and allow all nonconstant polynomials to factor completely into linear
factors over \(\mathbb{C}\).

Important additional distinctions include

\[
\boxed{
\text{rational versus irrational},
}
\]

\[
\boxed{
\text{algebraic versus transcendental},
}
\]

\[
\boxed{
\text{countable versus uncountable},
}
\]

and

\[
\boxed{
\text{discrete versus continuous}.
}
\]

Modular systems introduce finite arithmetic:

\[
\boxed{
\mathbb{Z}/n\mathbb{Z}.
}
\]

When \(p\) is prime,

\[
\boxed{
\mathbb{F}_p
}
\]

is a finite field.

The choice of number system determines which operations, equations, limits,
and algebraic constructions are available.

%%%%%%%%%%%%%%%%%%%%%%%%%%%%%%%%%%%%%%%%%%%%%%%%%%%%%%%%%%%%
\subsection{Section Connections}
\label{subsec:number-systems-connections}
%%%%%%%%%%%%%%%%%%%%%%%%%%%%%%%%%%%%%%%%%%%%%%%%%%%%%%%%%%%%

\chapterconnection{
Number Systems consolidates the numerical structures introduced throughout
the textbook, beginning with the foundational inclusion
\(\mathbb{N}\subset\mathbb{Z}\subset\mathbb{Q}\subset\mathbb{R}
\subset\mathbb{C}\). Integer and modular systems connect with Number Theory,
fields and extensions connect with Algebra, completeness and the real line
connect with Analysis, complex numbers connect with Complex Analysis,
finite fields connect with algebra and coding-related mathematics, and
floating-point representations connect with Numerical and Computational
Mathematics. The next reference section collects important Mathematical
Constants and their common values, definitions, and uses.
}

%%%%%%%%%%%%%%%%%%%%%%%%%%%%%%%%%%%%%%%%%%%%%%%%%%%%%%%%%%%%
% SECTION SOURCE TRACKING
%%%%%%%%%%%%%%%%%%%%%%%%%%%%%%%%%%%%%%%%%%%%%%%%%%%%%%%%%%%%

%------------------------------------------------------------
% 14.3 Number Systems
%------------------------------------------------------------
%
% Source basis:
%   - Standard elementary and advanced number-system conventions.
%   - Standard algebraic, analytic, complex, modular, and computational
%     number structures.
%
% Used for:
%   - natural numbers;
%   - whole numbers;
%   - integers;
%   - positive and negative integers;
%   - rational numbers;
%   - terminating and repeating decimals;
%   - irrational numbers;
%   - real numbers;
%   - completeness;
%   - density;
%   - algebraic numbers;
%   - transcendental numbers;
%   - complex numbers;
%   - conjugates and modulus;
%   - polar and exponential forms;
%   - roots of unity;
%   - Fundamental Theorem of Algebra;
%   - closure properties;
%   - fields and ordered fields;
%   - modular arithmetic;
%   - finite fields;
%   - extended real numbers;
%   - extended complex plane;
%   - Gaussian integers;
%   - quadratic extensions;
%   - countability;
%   - uncountability;
%   - cardinality;
%   - discrete and continuous number sets;
%   - vector spaces over number systems;
%   - matrix spaces;
%   - positional number systems;
%   - binary and hexadecimal notation;
%   - floating-point representation;
%   - even, odd, prime, and composite integers;
%   - Fundamental Theorem of Arithmetic;
%   - classification exercises.
%
% Notes:
%   - Mathematical Constants is developed next in Section 14.4.
%   - Algebraic Identities follows in Section 14.5.
%
%%%%%%%%%%%%%%%%%%%%%%%%%%%%%%%%%%%%%%%%%%%%%%%%%%%%%%%%%%%%